\section*{J.1 质量保证总则}
高精度参数确定系统的质量保证应覆盖模型、算法、数据和运行环境四个层面。每次发布前必须通过回归测试、稳定性测试和边界值测试，并形成可追溯记录。
\clearpage

\section*{J.2 数据质量核查}
建议为输入数据设置完整性、有效性和一致性三类规则。完整性确保关键字段不缺失；有效性确保数值范围合理；一致性确保跨通道时间戳同步。任一规则触发告警时，系统应自动降级处理。
\clearpage

\section*{J.3 模型质量核查}
模型核查包括结构正确性与参数合理性。结构正确性可通过单位一致性与极限行为分析验证；参数合理性可通过先验区间与历史分布比对验证。必要时触发人工复核流程。
\clearpage

\section*{J.4 算法质量核查}
算法质量核查重点关注收敛稳定性、异常恢复能力和计算效率。建议使用标准测试集按日执行自动化测试，并统计失败率、重启率和平均耗时的趋势变化，发现异常及时回滚。
\clearpage

\section*{J.5 运行时监控}
运行时应实时监控残差、参数漂移、阻尼系数与迭代次数。通过控制图和滑动窗口阈值检测可提前识别失稳趋势。对于持续异常，应自动切换到保守参数或冻结更新。
\clearpage

\section*{J.6 故障注入与演练}
建议定期开展故障注入演练，包括数据缺失、噪声激增、节点中断和配置错误等场景。通过演练验证告警链路与恢复机制，确保线上故障能够在可控时间内恢复。
\clearpage

\section*{J.7 版本管理与审计}
发布版本应绑定代码提交、配置文件和模型参数快照。审计记录需包含发布人、审批人、发布时间与变更摘要。出现异常时可快速定位到具体版本并执行回滚。
\clearpage

\section*{J.8 持续改进机制}
建立“问题复盘--策略更新--验证闭环”机制。对失败案例进行分类统计，持续优化初值策略、阻尼规则与数据预处理流程。通过周期性迭代，逐步提升系统鲁棒性与可维护性。
\clearpage

\section*{J.9 部署清单示例}
上线前清单可包括：输入通道检查、参数边界检查、日志落盘检查、监控告警检查、回滚脚本检查。清单化执行可减少人为遗漏，提升发布可靠性。
\clearpage

\section*{J.10 交付说明}
最终交付应包含论文文档、实验脚本、配置模板、运行手册和故障处理手册。对于关键结果，需提供复现实验步骤，确保第三方可在独立环境下复现主要结论。
